% ============================================================
% Formulario, argomento 10: variabili aleatorie continue (dispensa, capitolo 10)
% ============================================================
\argomento{Variabili aleatorie continue: densità e CDF}{Dispensa: capitolo 10}[Densità e CDF]

\voce{puntozero}{Un singolo valore ha probabilità zero}{variabile continua}
  {P(X = x) = 0 \quad \text{per ogni } x}
  {\variante{gli estremi non contano}{P(a_1 \leq X \leq a_2) = P(a_1 < X < a_2) = P(a_1 \leq X < a_2)}}

\voce{pdf}{Densità di probabilità (PDF)}{probabilità per unità di lunghezza}
  {f_X(x) = \lim_{\Delta x \to 0} \frac{P\left( x - \frac{\Delta x}{2} \leq X \leq x + \frac{\Delta x}{2} \right)}{\Delta x}}
  {\variante{intervallo piccolo}{P\left( x - \tfrac{\Delta x}{2} \leq X \leq x + \tfrac{\Delta x}{2} \right) \simeq
                                f_X(x)\, \Delta x}}

\voce{intervallo}{Probabilità di un intervallo}{l'area sotto la densità}
  {P(a_1 \leq X \leq a_2) = \int_{a_1}^{a_2} f_X(t) \dd t}
  {}

\voce{proprietapdf}{Proprietà della densità}{}
  {f_X(x) \geq 0 \qquad \int_{-\infty}^{\infty} f_X(t) \dd t = 1}
  {\variante{la densità può superare 1}{\text{per esempio } f_X(x) = 2 \ \text{su } \left[0, \tfrac{1}{2}\right]}}

\voce{cdf}{Funzione di distribuzione cumulativa (CDF)}{area a sinistra di $x$}
  {F_X(x) = P(X \leq x) = \int_{-\infty}^{x} f_X(t) \dd t \qquad f_X(x) = \frac{\dd}{\dd x} F_X(x)}
  {\variante{probabilità di un intervallo}{P(a_1 \leq X \leq a_2) = F_X(a_2) - F_X(a_1)}
   \variante{CCDF, area a destra di $x$}{\overline{F}_X(x) = P(X > x) = 1 - F_X(x), \quad
                                         f_X(x) = -\frac{\dd}{\dd x}\, \overline{F}_X(x)}}

\voce{proprietacdf}{Proprietà della CDF}{}
  {0 \leq F_X(x) \leq 1 \qquad F_X(-\infty) = 0 \qquad F_X(+\infty) = 1 \qquad x_1 < x_2 \Rightarrow F_X(x_1) \leq F_X(x_2)}
  {\variante{variabile continua}{F_X \text{ è continua}}
   \variante{variabile discreta}{F_X(x) = P(X \leq x) \text{ è una scala, con un salto alto } p_X(x) \text{ in ogni valore}}}

\voce{mediacontinua}{Media di una variabile continua}{se l'integrale converge}
  {E[X] = \mu_X = \int_{-\infty}^{\infty} x\, f_X(x) \dd x}
  {\variante{densità simmetrica attorno a $c$ (se la media esiste)}{E[X] = c}}
